\section{树递归神经网络的反思与展望}

\subsection{为什么树递归网络没有成为主流}

Manning 坦诚地分析了树递归神经网络最终未能竞争过 Transformer 的原因：

\begin{enumerate}
    \item \textbf{严格的上下文无关骨架}：信息只能沿树结构流动，没有像 Transformer 中 attention 那样的全局信息交互
    \item \textbf{依赖句法分析}：需要预先获得（或同时学习）句法树结构，增加了系统复杂度
    \item \textbf{难以并行化}：树结构的计算具有天然的序列依赖性
    \item \textbf{Transformer 的通用性}：Transformer 的 self-attention 在每一层都允许任意位置之间的信息交互，这种更灵活的信息流在实践中表现更好
\end{enumerate}

\begin{warningbox}{灵活性 vs 语言学正确性}
树递归网络在\textbf{语言学正确性}方面更有优势——它能更准确地建模否定、量词等语言现象，因为它尊重了语言的层次结构。但 Transformer 在\textbf{灵活性和规模化}方面具有压倒性优势。这揭示了一个深层矛盾：对于当前的深度学习范式，更灵活（但可能``不那么正确''）的模型往往比更有结构约束（但``更正确''）的模型表现更好。即使到今天，Transformer 模型在处理否定等现象时仍然不如树递归网络。
\end{warningbox}

\subsection{未来方向}

Manning 表示他仍然对一个问题感兴趣：能否将树结构的优势与 Transformer 的灵活性结合起来？这个方向的研究可能包括：

\begin{itemize}
    \item 在 Transformer 中引入结构化的归纳偏置
    \item 使用树结构作为 attention 的指导或正则化
    \item 结合组合语义的思想来增强大语言模型对否定、量词等现象的处理能力
\end{itemize}

\subsection{本章小结}

树递归神经网络因其严格的树结构约束和不够灵活的信息流而未能成为主流。然而，它在语言学正确性方面的优势（特别是否定建模）至今仍具有启发价值。将结构化的语言学知识与灵活的神经网络架构相结合，仍然是一个开放的研究方向。

%% ========== 总结与延伸 ==========

